\subsection{La dirección seleccionada por el registro dodecafásico}
\label{exc:k-direccion}

El registro \(K\) tiene una función geométrica que complementa su
evaluación racional y su lectura de soportes. Una vez construida la
representación de las doce posiciones, sus componentes determinan una
dirección en el espacio de variaciones de media nula. La reducción
\(12\to11\) elimina el modo uniforme; la reducción siguiente
\(11\to10\) elimina la dirección que selecciona el propio registro.
Las dos operaciones tienen, por tanto, antecedentes diferentes.

Para declarar completamente la coordenatización utilizada, consideremos el grupo
\(A_5\) de permutaciones pares de cinco objetos y el subgrupo
\(C_5=\langle(1\,2\,3\,4\,5)\rangle\). Las doce clases izquierdas
\(r_pC_5\) se ordenan mediante estos representantes, escritos como
listas de imágenes:
\[
\begin{array}{c|c@{\qquad}c|c}
p&r_p&p&r_p\\ \hline
1&(4,3,5,2,1)&7&(5,4,3,2,1)\\
2&(4,2,1,3,5)&8&(1,3,4,2,5)\\
3&(1,2,4,5,3)&9&(4,2,3,5,1)\\
4&(2,5,3,4,1)&10&(3,2,5,4,1)\\
5&(4,3,1,5,2)&11&(4,5,2,3,1)\\
6&(5,1,2,3,4)&12&(5,3,2,4,1).
\end{array}
\]
La acción por multiplicación izquierda define una representación
ortogonal \(\rho:A_5\to O(\RR^{12})\):
\(\rho(a)e_p=e_q\) cuando \(ar_pC_5=r_qC_5\).
Esta coordenatización de representación actúa sobre el registro ya obtenido;
los caracteres del grupo no se utilizan para elegir sus doce valores.

Sean \(\chi_3\) los valores del carácter tridimensional
\[
\begin{array}{c|ccccc}
\text{clase}&1&2^2&3&5_a&5_b\\ \hline
\chi_3&3&-1&0&(1+\sqrt5)/2&(1-\sqrt5)/2.
\end{array}
\]
La clase \(5_a\) contiene el generador de \(C_5\) mostrado y
\(5_b\) contiene su cuadrado. Esta convención distingue los dos
sectores tridimensionales conjugados de la representación. Definimos
\begin{equation}
 P_3=\frac3{60}\sum_{a\in A_5}\chi_3(a^{-1})\rho(a),
 \qquad P_{11}=I_{12}-\frac1{12}J_{12}.
 \label{eq:k-proyectores-incidencia}
\end{equation}
Ambas matrices quedan determinadas por sumas finitas en
\(\QQ(\sqrt5)\). La aparición de este cuerpo cuadrático pertenece
a la realización de la representación, posterior a la generación
regional de \(\varphi\).

\begin{lemma}[Proyección del registro sobre el sector tridimensional]
\label{lem:k-sector-no-nulo}
Se tienen
\[
 P_3=P_3^{\mathsf T}=P_3^2,\quad
 P_3\mathbf1=0,\quad \operatorname{tr}P_3=3,
 \qquad P_3P_{11}=P_{11}P_3=P_3.
\]
Para el registro \eqref{exc:k}, el vector
\begin{equation}
 u_K=P_3P_{11}K
 \label{eq:k-vector-direccional}
\end{equation}
es no nulo. Su norma exacta es
\begin{equation}
 \lVert u_K\rVert^2
 =\frac{6638585+2275584\sqrt5}{20}>0.
 \label{eq:k-norma-direccion}
\end{equation}
\end{lemma}
\begin{proof}
La suma de carácter en \eqref{eq:k-proyectores-incidencia} es el
proyector isótipo. El carácter real y la sustitución \(a\mapsto a^{-1}\)
dan la simetría. La multiplicidad del sector en la acción sobre
\(A_5/C_5\) es la dimensión de sus vectores fijos por \(C_5\),
\[
 \frac15\left(3+2\frac{1+\sqrt5}{2}
                 +2\frac{1-\sqrt5}{2}\right)=1.
\]
Por ello el rango es tres. Su ortogonalidad al carácter trivial
da \(P_3\mathbf1=0\) y las dos composiciones con \(P_{11}\).
También pueden verificarse estas identidades multiplicando las
matrices de la suma finita.

Finalmente, \(\sum K_p=6263\) y la sustitución de las doce
componentes en esa misma suma da
\[
 \lVert u_K\rVert^2=K^{\mathsf T}P_3K
 =\frac1{20}\sum_{a\in A_5}
                \chi_3(a^{-1})K^{\mathsf T}\rho(a)K
 =\frac{6638585+2275584\sqrt5}{20}.
\]
Esta es una evaluación exacta de sesenta términos especificados,
reproducida por el certificado adjunto con aritmética racional
cuadrática. Los dos coeficientes positivos prueban la no nulidad.
\end{proof}

\begin{definition}[Dirección dodecafásica y complemento transversal]
La dirección seleccionada por \(K\) en la coordenatización declarada es
\(\ell_K=\RR u_K\). Su complemento transversal dentro de
\(V_{11}=\mathbf1^\perp\) es
\[
 V_{10}(K)=V_{11}\cap u_K^\perp.
\]
\end{definition}

\begin{proposition}[Reducción ortogonal determinada por \(K\)]
\label{prop:k-reduccion-diez}
El operador
\begin{equation}
 P_{10}(K)=P_{11}
       -\frac{u_Ku_K^{\mathsf T}}{\lVert u_K\rVert^2}
 \label{eq:k-proyector-diez}
\end{equation}
es el proyector ortogonal de rango diez sobre \(V_{10}(K)\). Se cumple
\[
 V_{11}=\ell_K\mathbin{\oplus^{\perp}}V_{10}(K),
 \qquad P_{10}P_{11}=P_{10},\quad P_{10}u_K=0.
\]
\end{proposition}
\begin{proof}
Sea \(E_K=u_Ku_K^{\mathsf T}/\lVert u_K\rVert^2\). La
no nulidad demostrada garantiza que está definido. Es simétrico,
\(E_K^2=E_K\), tiene rango uno y, como \(u_K\in V_{11}\),
\(P_{11}E_K=E_KP_{11}=E_K\). De aquí
\[
 (P_{11}-E_K)^2=P_{11}-2E_K+E_K=P_{11}-E_K.
\]
Su imagen es \(V_{11}\cap u_K^\perp\) y su traza es
\(11-1=10\). La descomposición y las identidades restantes se
obtienen por proyección sobre los dos sumandos ortogonales.
\end{proof}

Este resultado precisa una función de \(K\) que su denominación como
registro no agotaba: además de representar reversiblemente datos
orientados y de seleccionar la tétrada \(B_K\), determina una
polarización lineal. La dirección depende del registro completo a
través de \(P_3K\), no únicamente de los cuatro puntos de su soporte.
Se mantiene así el orden registro, representación, dirección y
reducción dimensional: \eqref{eq:k-vector-direccional} define la
dirección y la proposición~\ref{prop:k-reduccion-diez} demuestra
la reducción y caracteriza su núcleo.

La transformación tiene una ley de covariancia precisa. Si se cambia
simultáneamente la coordenatización por una matriz ortogonal de permutación \(S\),
\(K'=SK\) y \(P'_3=SP_3S^{-1}\), entonces
\[
 u_{K'}'=Su_K,\qquad P'_{10}(K')=SP_{10}(K)S^{-1}.
\]
Por otra parte, sustituir \(u_K\) por \(a u_K\), con \(a\ne0\),
no cambia \(E_K\). La dirección retiene una polarización y no
permite recuperar por sí sola ni la amplitud de \(u_K\) ni el
registro íntegro. La reversibilidad de \(\kappa\leftrightarrow K\)
permanece en su dominio propio.

La realización geométrica aquí establecida es una reducción de espacios
con producto interno. Su empleo posterior como dirección compacta
requiere transportar esta recta a la realización física correspondiente.
El apartado \ref{exc:pantallas-dualidad} expone el mapa algebraico
que mantiene coordinadas esta reducción y la pantalla reticular.
